\chapter{Construire une estimation probabiliste de la gravité}
\label{chap:protocole-ml}

Les analyses précédentes comparent les groupes séparément. Elles montrent des écarts,
mais ne disent pas comment plusieurs caractéristiques, considérées ensemble, permettent
d'estimer une issue individuelle. Le volet prédictif répond à cette seconde question.
Son résultat est une probabilité de décès ou d'hospitalisation parmi les usagers
d'accidents corporels enregistrés. Aucun scénario de prédiction avant un trajet n'est
défini dans cette étude.

\section{Ce que le modèle cherche à estimer}
Pour un usager décrit par $x$, la quantité estimée est
\[
  p(x)=\mathbb{P}(Y=1\mid X=x,\ \text{usager présent dans BAAC}).
\]
La cible $Y$ conserve la définition du chapitre~\ref{chap:cadrage}. Les gravités
inconnues sont exclues de l'ajustement et des scores, après comptage. Le modèle ne
distingue donc pas les décès des hospitalisations à l'intérieur de la classe positive.

La comparaison est volontairement limitée à une référence constante et à une
régression logistique. Cette dernière combine plusieurs caractéristiques, produit
directement des probabilités et permet une lecture des coefficients. Un modèle plus
complexe n'est pas nécessaire pour vérifier si ces informations apportent déjà un gain.
Ce choix limite aussi le nombre d'expériences utilisées pour sélectionner une solution.

\section{Séparer les années avant l'apprentissage}
L'apprentissage porte sur 2019-2022, la validation sur 2023 et le test final sur 2024.
La validation sert à choisir le réglage du modèle. Le test sert ensuite à mesurer les
performances de cette configuration sur une année ultérieure.
\begin{table}[htbp]
\centering\small
\renewcommand{\arraystretch}{1.18}
\caption{Populations utilisées pour la modélisation, après exclusion des gravités inconnues.}\label{tab:ml-population}
\begin{tabular*}{\textwidth}{@{\extracolsep{\fill}}lrrrr@{}}
\toprule
Période & Cibles connues & Exclues & Cas positifs & Part (\%) \\
\midrule
2019-2022 & 467\,079 & 269 & 83\,081 & 17,79 \\
2023 & 118\,101 & 96 & 21\,218 & 17,97 \\
2024 & 117\,278 & 0 & 21\,024 & 17,93 \\
\bottomrule
\end{tabular*}
\sourcecalcul{notebooks 06 et 07 ; les exclusions portent uniquement sur la cible inconnue}
\end{table}


Une séparation par année maintient tous les usagers d'un accident dans la même
partition. Elle évite qu'un conducteur soit utilisé pour apprendre tandis qu'un
passager du même accident servirait à évaluer le modèle. Elle conserve également les
changements entre millésimes dans l'évaluation.

L'année 2024 a toutefois été examinée lors des audits de structure. Ce test est donc
\textbf{rétrospectif et chronologique}, et non une expérience prospective sur un jeu
resté totalement inconnu. Ses scores n'ont pas servi à sélectionner les variables,
les valeurs de régularisation ou le seuil. Maintenant qu'ils sont consultés, cette
année ne peut plus être réutilisée comme un nouveau test vierge.

\section{Choisir des variables et rendre leur représentation explicite}
Onze variables sources sont retenues avant la comparaison des scores. Elles décrivent
l'usager, le véhicule associé et les circonstances de l'accident.

\begin{table}[htbp]\centering\small
\caption{Variables du premier modèle et représentation retenue.}\label{tab:variables-ml}
\renewcommand{\arraystretch}{1.17}
\begin{tabularx}{\textwidth}{@{}lLL@{}}\toprule
Famille & Variables sources & Représentation et précautions\\\midrule
Usager & \code{age_approx}, \code{sexe}, \code{catu} & Âge en huit tranches fixes ; sexe et catégorie comme modalités.\\
Véhicule & \code{catv} & Catégorie du véhicule associé ; véhicule heurtant pour un piéton.\\
Contexte & \code{lum}, \code{agg}, \code{int}, \code{atm} & Lumière, agglomération, intersection et conditions atmosphériques.\\
Temps & \code{mois_accident}, \code{jour_semaine}, \code{heure_accident} & Mois, jour et heure traités comme catégories.\\\bottomrule
\end{tabularx}
\sourcecalcul{notebook 06 et module \code{src/baac_ml.py}}
\end{table}

L'âge reprend les groupes de l'exploration~: 0-17, 18-24, 25-34, 35-44, 45-54,
55-64, 65-74 et 75 ans ou plus. Les bornes ne sont pas optimisées sur 2023.
Cette représentation rend les effets lisibles, au prix d'une perte de précision
à l'intérieur de chaque tranche. Le traitement catégoriel du temps évite d'imposer
une relation linéaire entre le numéro d'une heure et les log-cotes de gravité.

Une valeur absente devient la modalité «~Non renseigné~». Il n'y a pas d'imputation
par une moyenne calculée sur les six années. L'encodage indicateur, ou \textit{one-hot},
crée une colonne 0/1 par modalité observée dans l'apprentissage. Avec toutes les
modalités conservées, la matrice comporte \textbf{117 colonnes}. Une catégorie nouvelle
est comptée et encodée en zéros pour son champ, sans réajuster l'encodeur.

Les identifiants, la provenance, l'année, la gravité et ses dérivés sont exclus.
Les indicateurs de qualité et variantes brutes déjà représentées ne sont pas ajoutés.
Les champs de lieux sont laissés hors de ce premier modèle en raison des difficultés
de représentation mises en évidence en 2023. Les équipements, manœuvres et détails du
choc sont également hors périmètre, sans être tous qualifiés de fuites de cible.

Ces variables proviennent de fiches décrivant un accident déjà enregistré. Leur
présence dans un fichier final ne prouve pas leur disponibilité à un instant de décision.
La représentation des valeurs manquantes peut notamment refléter la saisie après
l'accident. Cette limite concerne l'interprétation et toute utilisation envisagée.

\section{Une référence constante et une régression logistique}
La référence attribue à tous les usagers la fréquence positive de l'apprentissage~:
\[
  \widehat p_0=\frac{83\,081}{467\,079}\simeq 0{,}17787.
\]
Elle ne connaît ni les fréquences de 2023 et 2024, ni les caractéristiques des usagers
évalués. Elle fournit un point de comparaison plus informatif qu'un score de modèle isolé.

La régression logistique transforme une combinaison des colonnes encodées $z_j(x)$
en probabilité~\cite{sklearn-logistique}~:
\[
  \widehat p(x)=\frac{1}{1+\exp[-\eta(x)]},\qquad
  \eta(x)=\beta_0+\sum_{j=1}^{117}\beta_j z_j(x).
\]
Le critère ajusté est la somme des pertes logarithmiques, avec une pénalisation L2
des coefficients. Sans pondération des observations, on peut l'écrire~:
\[
  \min_{\beta_0,\beta}\left\{
  \sum_{i=1}^{n}\big[-y_i\log\widehat p_i-(1-y_i)\log(1-\widehat p_i)\big]
  +\frac{1}{2C}\sum_{j=1}^{117}\beta_j^2\right\}.
\]
Un petit $C$ impose une pénalisation plus forte. Chaque usager contribue au critère~;
un accident avec plusieurs usagers contribue donc davantage qu'un accident avec un seul
usager. Cette pondération correspond à l'unité retenue, mais ne rend pas les observations
d'un même accident indépendantes.

L'encodeur et le modèle sont réunis dans une pipeline ajustée uniquement sur 2019-2022.
Le solveur \code{lbfgs} utilise une tolérance de $10^{-6}$ et au plus 1\,500 itérations.
Une non-convergence interrompt la comparaison. Aucune pondération ou duplication des
classes n'est appliquée~: aucun coût métier des erreurs n'a été défini.

\section{Évaluer des probabilités et un classement}
La log-loss est le critère principal de sélection~\cite{sklearn-logloss}~:
\[
 \mathrm{LL}=-\frac1n\sum_{i=1}^n
 \left[y_i\log\widehat p_i+(1-y_i)\log(1-\widehat p_i)\right].
\]
Elle pénalise particulièrement une probabilité très assurée lorsqu'elle contredit
l'issue observée. Le score de Brier mesure l'écart quadratique moyen entre probabilités
et cibles~\cite{sklearn-brier}~:
\[
 \mathrm{Brier}=\frac1n\sum_{i=1}^n(\widehat p_i-y_i)^2.
\]
Pour ces deux indicateurs, une valeur plus petite est préférable. Leur amélioration ne
suffit toutefois pas à établir une bonne calibration dans chaque groupe.

Deux indicateurs complètent la lecture. La ROC-AUC mesure la capacité à classer un cas
positif au-dessus d'un cas négatif, en tenant compte des ex æquo~\cite{sklearn-roc}.
L'\textit{average precision} (AP) résume la courbe précision-rappel par une somme
pondérée par les gains de rappel~\cite{sklearn-ap}. Elle n'est pas une aire calculée
par interpolation trapézoïdale. Pour un score constant, elle vaut ici la proportion
positive de l'année évaluée, même si la probabilité constante vient de l'apprentissage.

\FloatBarrier
\bilan{L'expérience est définie avant le test~: une cible binaire, onze variables,
un encodeur appris sur les années antérieures, une référence constante et une seule
famille de modèles. Elle évalue un apport prédictif rétrospectif, avec des limites de
collecte et de disponibilité des variables explicites.}
